\documentclass[11pt,a4paper]{article}
\usepackage{fullpage}
\usepackage[utf8]{inputenc}
\usepackage[T1]{fontenc}
\usepackage[french]{babel} \frenchbsetup{og=«,fg=»}
\usepackage{graphicx}
\usepackage{amsthm}
\usepackage{amsmath}
\usepackage{amssymb}
\usepackage{stmaryrd}
\usepackage{proof}
\usepackage{booktabs}


\theoremstyle{definition}
\newtheorem{definition}{Définition}[section]
\newtheorem{theorem}{Théorème}[section]
\newtheorem{lemma}[theorem]{Lemme}



\title{Étude d'une mesure d'espace pour le lambda calcul}
\author{Ewen DUFOUR\\[1ex]
{\small Encadré par}\\ {\small Thibaut BALABONSKI}}
\date{Stage de M1 MPRI\\mars--juillet 2025}
\begin{document}

\maketitle

\begin{abstract}
  Nous sommes beaux Quelque chose de bien orthographiée.
\end{abstract}

\tableofcontents
\section{Introduction}
\section{Le Linear Substitution Calculus}

\begin{definition}[Linear Substitution Calculus]
Les termes du LSC sont donnés par la grammaire~:
\begin{align*}
    t ::= x~|~t~t~|~\lambda x.t~|~t[x / t]
\end{align*}
Les règles de réduction sont~:
\begin{center}
\begin{tabular}{l l l l }
$(\lambda x.t)Lu$ &$\mapsto_{d\beta}$ &$t[x / u]L$ & \\
$\mathcal{C}\llbracket x \rrbracket [x / u]$ &$\mapsto_{ls}$ &$\mathcal{C}\llbracket u \rrbracket [x / u]$ & \\
$t[x / u]$ &$\mapsto_{gc}$ &$t$ & if $x\notin fv(t)$
\end{tabular}
\end{center}
 où $L$ est une liste de substitutions possiblement vide, $\mathcal{C}$ est un contexte et $fv(t)$ dénote les variables libres du terme $t$.
\end{definition}
\begin{definition}[Free variables]
\begin{align*}
    fv(x) &= \{ x \} \\
    fv(\lambda x.t) &= fv(t) \setminus \{ x \} \\
    fv(t~u) &= fv(t) \cup fv(u) \\
    fv(t[x/u]) &= (fv(t) \setminus \{ x \}) \cup fv(u) \\
\end{align*}
\end{definition}

\section{Le Weak Call By Value LSC}
Nous considérons les termes suivant comme des valeurs~:
\begin{align*}
    v ::= (\lambda x.t)~|~v[x/t]
\end{align*}
\begin{definition}[Contextes de réductions]
Nous définissons les contextes de réduction ainsi~:
\begin{align*}
    \mathcal{C} ::= \square~|~\mathcal{C}~t~|~v~\mathcal{C}~|~\mathcal{C}[x/t]
\end{align*}
Ou de manière équivalente :
\begin{align*}
    \mathcal{C} ::= \square~|~\mathcal{C}[\square~t]~|~\mathcal{C}[v~\square]~|~\mathcal{C}[\square[x/t]]
\end{align*}
\end{definition}
\begin{figure}[!h]
    \centering
    \begin{align*}
    &\infer[LA]{t~v\rightarrow t'~u}{t \rightarrow t'}
 &\infer[RA]{v~t\rightarrow v~t'}{t \rightarrow t'}\\
&\infer[GC]{t[x/u]\rightarrow t'}{t \rightarrow t' & x \notin fv(t') }
&\infer[US]{t[x/u]\rightarrow t'[x/u]}{t \rightarrow t' & x \in fv(t')}\\
&\infer[LS]{\mathcal{C}\llbracket x \rrbracket [x/u]\rightarrow \mathcal{C}\llbracket u \rrbracket [x/u]}{x \in fv(\mathcal{C}\llbracket u \rrbracket) }
&\infer[LSGC]{\mathcal{C}\llbracket x \rrbracket [x/u]\rightarrow \mathcal{C}\llbracket u \rrbracket}{x \notin fv(\mathcal{C}\llbracket u \rrbracket) }\\
&\infer[DB]{(\lambda x.t)L~v \rightarrow t[x/v]L}{x\in fv(t)}
&\infer[DBGC]{(\lambda x.t)L~v \rightarrow tL}{x\notin fv(t)}
\end{align*}
    \caption{Règles de réduction du WCBVLSC}
\end{figure}

\begin{definition}[Propreté]
\begin{align*}
    cl(x) &= true\\
    cl(\lambda x.t) &= cl(t)\\
    cl(t~u) &= cl(t) \wedge cl(u)\\
    cl(t[x/u]) &= cl(t) \wedge x\in fv(t) \wedge cl(u)
\end{align*}
Un terme est dit propre si et seulement si $cl(t)=true$.
\end{definition}

\begin{lemma}[Plongeon]
  Soient $\mathcal{C}$, $\mathcal{D}$ des contextes et $t$ un terme.
\[\langle -,\mathcal{C}[t], \mathcal{D}\rangle_0 \rightsquigarrow^* \langle -,t,\mathcal{C} :: \mathcal{D}\rangle_0\]
\end{lemma}
\begin{proof}
  Par induction sur la forme de $\mathcal{C}$.
  \begin{itemize}
    \item \underbar{Cas $\mathcal{C} = \square$~:} immédiat.\\
    \item \underbar{Cas $\mathcal{C} = \mathcal{C}_1[x/v]$~:}\\
  \begin{align*}
    \langle -, \mathcal{C}[t], \mathcal{D}\rangle_0 = &\langle -, \mathcal{C}_1[t][x/v], \mathcal{D}\rangle_0 \\
    \rightsquigarrow_{us} &\langle -, \mathcal{C}_1[t], \square[x/v]:: \mathcal{D}\rangle_0 \\
    \rightsquigarrow^* &\langle -, t, \mathcal{C}_1:: \square[x/v]:: \mathcal{D}\rangle_0 \quad \text{Par hypothèse d'induction.}
  \end{align*}
    \item \underbar{Cas $\mathcal{C} = \mathcal{C}_1~u$~:}\\
  \begin{align*}
    \langle -, \mathcal{C}[t], \mathcal{D}\rangle_0 = &\langle -, \mathcal{C}_1[t]~u, \mathcal{D}\rangle_0 \\
    \rightsquigarrow_{lapp} &\langle -, \mathcal{C}_1[t], \square~u:: \mathcal{D}\rangle_0 \\
    \rightsquigarrow^* &\langle -, t, \mathcal{C}_1:: \square~u:: \mathcal{D}\rangle_0 \quad \text{Par hypothèse d'induction.}
  \end{align*}
    \item \underbar{Cas $\mathcal{C} = v~\mathcal{C}_1$~:}\\
  \begin{align*}
    \langle -, \mathcal{C}[t], \mathcal{D}\rangle_0 = &\langle -, v~\mathcal{C}_1[t], \mathcal{D}\rangle_0 \\
    \rightsquigarrow_{lapp} &\langle -, \mathcal{C}_1[t], v~\square:: \mathcal{D}\rangle_0 \\
    \rightsquigarrow^* &\langle -, t, \mathcal{C}_1:: v~\square:: \mathcal{D}\rangle_0 \quad \text{Par hypothèse d'induction.}
  \end{align*}
  \end{itemize}
\end{proof}

\begin{lemma}[Reconstruction]
  Soient $\mathcal{C}$ et $\mathcal{D}$ des contextes et $t$ un terme.
  \[\mathcal{C}[t]~\text{propre}~\Rightarrow \langle-,t,\mathcal{C}::\mathcal{D}\rangle_2 \rightsquigarrow^* \langle-,\mathcal{C}[t],\mathcal{D}\rangle_2\]
\end{lemma}
\begin{proof}
  Par induction sur la forme de $\mathcal{C}$.
  \begin{itemize}
    \item \underbar{Cas $\mathcal{C} = \square$~:} immédiat.
    \item \underbar{Cas $\mathcal{C} = \mathcal{C}_1[x/v]$~:}\\
      Nous avons~:
      \begin{align*}
      \mathcal{C}[t]~\text{propre} &\Leftrightarrow \mathcal{C}_1[t][x/v]~\text{propre}\\
      &\Leftrightarrow x\in fv(\mathcal{C}_1[t]) \wedge \mathcal{C}_1[t]~\text{propre}
      \end{align*}
      Et par conséquent~:
      \begin{align*}
        \langle -, t, \mathcal{C} :: \mathcal{D}\rangle_2 = &\langle -, t, \mathcal{C}_1[x/v]:: \mathcal{D}\rangle_2 \\
        = &\langle -, t, \mathcal{C}_1:: \square[x/v] :: \mathcal{D}\rangle_2 \\
        \rightsquigarrow^* &\langle -, \mathcal{C}_1[t], \square[x/v]:: \mathcal{D}\rangle_2 \quad &\text{Par hypothèse d'induction.}\\
        \rightsquigarrow_{bsubst} &\langle -, \mathcal{C}_1[t][x/v], \mathcal{D}\rangle_2 \quad &\text{Car}~x\in fv(\mathcal{C}_1[t])
      \end{align*}
    \item \underbar{Cas $\mathcal{C} = \mathcal{C}_1~u$~:}\\
      Nous avons~:
      \begin{align*}
        \mathcal{C}[t]~\text{propre} &\Leftrightarrow \mathcal{C}_1[t]~u~\text{propre}\\
        &\Rightarrow \mathcal{C}_1[t]~\text{propre}
      \end{align*}
      Ce qui nous donne~:
      \begin{align*}
        \langle -, t, \mathcal{C} :: \mathcal{D}\rangle_2 = &\langle -, t, \mathcal{C}_1~u:: \mathcal{D}\rangle_2 \\
        = &\langle -, t, \mathcal{C}_1:: \square~u :: \mathcal{D}\rangle_2 \\
        \rightsquigarrow^* &\langle -, \mathcal{C}_1[t], \square~u:: \mathcal{D}\rangle_2 \quad &\text{Par hypothèse d'induction.}\\
        \rightsquigarrow_{blapp} &\langle -, \mathcal{C}_1[t]~u, \mathcal{D}\rangle_2 \quad &\text{Par propreté de}~\mathcal{C}
      \end{align*}
    \item \underbar{Cas $\mathcal{C} = v~\mathcal{C}_1$~:} Analogue au cas précédent.
  \end{itemize}
\end{proof}
\begin{lemma}[Préservation des substitutions]
  Soient $\mathcal{C}$ un contexte, $x$ et $y$ des variables et t un terme.
  \[
    y \neq x~\wedge~y\in fv(\mathcal{C}[x]) \Rightarrow y \in fv(\mathcal{C}[t])
  \]
\end{lemma}
\begin{proof}
Par induction sur la forme de $\mathcal{C}$
 \begin{itemize}
    \item \underbar{Cas $\mathcal{C} = \square$~:} Prémisse fausse.
    \item \underbar{Cas $\mathcal{C} = \mathcal{C}_1[z/v]$~:}\\
      Comme $y\in fv(\mathcal{C}[x])$, nous avons deux cas.
      \begin{itemize}
        \item \underbar{Cas $y \in fv(v)$~:}\\ 
          En particulier, $y \in fv(\mathcal{C}_1[t][z/v]) \Leftrightarrow y \in fv(\mathcal{C}[t])$
        \item \underbar{Cas $y\in fv(\mathcal{C}_1[x]) \wedge y \neq z$~:}\\
          Par hypothèse d'induction, $y\in fv(\mathcal{C}_1[t])$.\\
          Comme $y \neq z$, $y\in fv(\mathcal{C}_1[t][z/v]) \Leftrightarrow y \in fv(\mathcal{C}[t]$ 
      \end{itemize}
    \item \underbar{Cas $\mathcal{C} = \mathcal{C}_1~u$~:}\\
      De nouveau, nous avons deux cas.
      \begin{itemize}
        \item \underbar{Cas $y \in fv(u)$~:}\\
          $ y \in fv(\mathcal{C}[t]~u) \Leftrightarrow y \in fv(\mathcal{C}[t])$
        \item \underbar{Cas $y \in fv(\mathcal{C}_1[x])$~:}
          Par hypothèse d'induction, $y \in fv(\mathcal{C}_1[t])$\\ 
          Donc, $y \in fv(\mathcal{C}_1[t]~u) \Leftrightarrow y \in fv(\mathcal{C}[t])$
      \end{itemize}
    \item \underbar{Cas $\mathcal{C} = v~\mathcal{C}_1$~:} Analogue au cas précédent.

  \end{itemize} 
\end{proof}

\begin{lemma}[Propreté des substitutions]
 Soient $\mathcal{C}$ un contexte, $t$ un terme et $x$ une variable.
\[\mathcal{C}\llbracket x \rrbracket~\text{propre}~\wedge~t~\text{propre}~\Rightarrow \mathcal{C}[t]~\text{propre}\]
\end{lemma}
\begin{proof}
  Par induction sur la forme de $\mathcal{C}$.
  \begin{itemize}
    \item \underbar{Cas $\mathcal{C} = \square$~:} trivial.
    \item \underbar{Cas $\mathcal{C} = \mathcal{C}_1[y/v]$~:}\\
      Comme $\mathcal{C}\llbracket x\rrbracket$ est propre, $\mathcal{C}_1\llbracket x \rrbracket$ et $v$ sont propres et $y\in fv(\mathcal{C}_1\llbracket x\rrbracket)$.
      En particulier, $x \neq y$, car $x$ est supposé libre dans $\mathcal{C}\llbracket x \rrbracket$.
      Par conséquence du lemme de préservation des substitutions, $y \in fv(\mathcal{C}_1[t])$.\\
      Par hypothèse d'induction, nous avons $\mathcal{C}_1[t]$ propre.\\
      Donc $\mathcal{C}[t]$ est propre.
    \item \underbar{Cas $\mathcal{C} = \mathcal{C}_1~u$~:}\\
      Nous avons~:
      \begin{align*}
        \mathcal{C}[t]~\text{propre} &\Leftrightarrow \mathcal{C}_1[t]~u~\text{propre}\\
        &\Rightarrow \mathcal{C}_1[t]~\text{propre}
      \end{align*}
      Ce qui nous donne~:
      \begin{align*}
        OOPS
      \end{align*}
    \item \underbar{Cas $\mathcal{C} = v~\mathcal{C}_1$~:} Analogue au cas précédent.
  \end{itemize}
\end{proof}

\section{Le Space Calculus}

\begin{definition}[Termes du Space Calculus]
Le Space calcul est défini par les termes donnés par la grammaire suivante~:
\begin{align*}
    l ::&= x \\
    &|~\lambda x.a \\
    &|~a~a\\
    a ::&= l^{\alpha}\\
    c ::&= a\\
        &|~c[x/c]\\
    t ::&= c\\
        &|~t~c
\end{align*}
ainsi que par les règles de réductions~:
\begin{align*}
&\infer[sea_v]{((t^\beta~x^\gamma)^\alpha e) \rightarrow (t^\beta e_{|t})~e(x)}{}
&\infer[sea_{\neg v}]{((t^\beta~u^\gamma)^\alpha e) \rightarrow (t^\beta e_{|t})~(u^\gamma e_{|u}) }{}\\\\
&\infer[\beta_{\neg w}]{((\lambda x .t^\beta)^\alpha e_t)~(u^\gamma e_u ) \rightarrow (t^\beta[x/u^\gamma e_u]e_t)}{x \in fv(t) }
&\infer[\beta_w]{((\lambda x .t^\beta)^\alpha e_t)~(u^\gamma e_u ) \rightarrow (t^\beta e_t)}{x \notin fv(t) }\\\\
&\infer[sub]{x^\alpha[x/u^\beta e] \rightarrow u^\beta e}{}
&\infer[cont]{\mathcal{C}[t] \rightarrow \mathcal{C}[t']}{t\rightarrow t'}
\end{align*}
où $e_{|t}$ dénote la restriction de l'environnement $e$ aux variables libres de $t$ et 
$\mathcal{C}$ est un contexte de réduction de la forme~:
\begin{align*}
  \mathcal{C} ::= \square~|~\mathcal{C}~c.
\end{align*}

\end{definition}
\begin{definition}[Restriction d'un environnement]
\begin{align*}
a_{|t} &= a\\
    (e[x/c])_{|t} &= e_{|t}[x/c] \quad &\text{si } x\in fv(t)\\
    (e[x/c])_{|t} &= e_{|t} \quad &\text{si } x\notin fv(t)
\end{align*}
\end{definition}

\section*{Annexe}

\begin{figure}[ht]
\centering
\resizebox{\textwidth}{!}{%
\begin{tabular}{@{}l l l l l l l l l l@{}}
\toprule
\textbf{arg} & \textbf{terme} & \textbf{contexte} & \textbf{état} & \textbf{transition}& \textbf{arg} & \textbf{terme} & \textbf{contexte} & \textbf{état} \\
\midrule

– & $t[ u ] $ & $C$ & 0 & $\rightarrow_{us}$ & – & $t$ & $\textsf{LS}(u) :: C$ & 0  \\

- & $t~u$ & $C$ & 0 & $\rightarrow_{lapp}$ & – & $t$ & $\textsf{RA}(u) :: C$ & 0  \\
- & $v~t$ & $C$ & 0 & $\rightarrow_{rapp}$ & – & $t$ & $\textsf{LA}(v) :: C$ & 0  \\
- & $n$  & $C$ & 0 & $\rightarrow_{ls}$ & - & $\uparrow^{n+1} t$ & $C$ & 2 & $\textsf{LS}(t)=\textsf{nth\_subst}(C)$ \\
- & $v_1~v_2$  & $C$ & 0 & $\rightarrow_{db}$ & $(v_2,0)$ & $v_1$ & $C$ & 1  \\

$(v_2, n)$ & $v_1[ t ]$ &$C$ & 1 & $\rightarrow_{rdb}$ & $(v_2, n+1)$ & $v_1$ & $\textsf{LS}(t) :: C$ & 1\\ 
$(v, n)$ & $\lambda.t $ &$C$ & 1 & $\rightarrow_{fdb}$ & - & $t$ & $\textsf{LS}(\uparrow^{n}v) :: C$ & 2\\

-&$v$ & $[]$ & 2 & $\rightarrow_{stop}$ & -& $v$ &$[]$ & 3 \\
-&$t$ & $[]$ & 2 & $\rightarrow_{dive}$ & -& $t$ &$[]$ & 0 \\
-&$t$ & $\textsf{LS}(u) :: C$ & 2 & $\rightarrow_{gc}$ & -& $\downarrow t$ &$C$ & 2 & $0\notin \textsf{fv}(t)$ \\
-&$t$ & $\textsf{LS}(u) :: C$ & 2 & $\rightarrow_{bsubst}$ & -& $t[ u ]$ &$C$ & 2 & $0\in \textsf{fv}(t)$ \\
-&$t$ & $\textsf{RA}(u) :: C$ & 2 & $\rightarrow_{blapp}$ & -& $t~u$ &$C$ & 2 \\
-&$t$ & $\textsf{LA}(u) :: C$ & 2 & $\rightarrow_{brapp}$ & -& $u~t$ &$C$ & 2 \\

\bottomrule
\end{tabular}
}
\caption{Transitions de la FLAME.}
\end{figure}
\end{document}

